\section{Arquitetura Alvo e Fluxo por Aba do Sistema}
\label{sec:arquitetura_alvo}

% ==============================================================================
% 1. ETAPA 1: O QUE É O FLUXO POR ABA NA ARQUITETURA ALVO
% ==============================================================================
\subsection{Etapa 1: O que é o ``Fluxo por Aba'' na Arquitetura Alvo?}

Dentro da metodologia TDSP (\textit{Team Data Science Process}), o ``Fluxo por Aba'' representa o desenho formal da engenharia de dados e software subjacente a cada funcionalidade, delimitando a jornada integral percorrida pela informação e detalhando a entrada, o processamento e a saída do sistema.

O objetivo dessa etapa reside em assegurar a sustentabilidade, auditabilidade e manutenibilidade do sistema (\textit{legacy system management}). O detalhamento garante que qualquer cientista de dados, engenheiro de Machine Learning, compreenda de imediato a finalidade de cada aba e como os dados trafegam pela arquitetura.

% ==============================================================================
% 2. ETAPA 2: O QUE CADA FLUXO PRECISA CONTER (ESTRUTURA PADRÃO)
% ==============================================================================
\subsection{Etapa 2: O que cada Fluxo precisa conter? (Estrutura Padrão)}

Para assegurar consistência metodológica, homogeneidade na documentação e facilidade de sustentação técnica, cada aba documentada na Arquitetura Alvo obedece a 5 componentes obrigatórios:

\begin{enumerate}[label=\textbf{\arabic*.}]
    \item \textbf{Objetivo \& User Story Vinculada:} Define com precisão o problema de negócio resolvido pela aba e qual persona destinatária ela atende (\textit{Candidato}, \textit{Recrutador} ou \textit{Professora}).
    \item \textbf{Fontes de Dados \& Artefatos:} Identifica formalmente a origem dos insumos consumidos, distinguindo arquivos brutos (\texttt{.csv}), dados pré-processados de alta performance (\texttt{.parquet}) e modelos serializados e metadados (\texttt{.pkl}).
    \item \textbf{Cadeia de Processamento (Pipeline de Lógica):} Detalha o conjunto de transformações matemáticas e operações vetoriais executadas (como ponderação léxica, vetorização TF-IDF, similaridade por cosseno, decomposição SVD com produto interno de fatores latentes, fusão empírica e ordenações).
    \item \textbf{Comportamento sob Limitações / Fallback:} Estipula como a aba reage em cenários de exceção, inicialização a frio (\textit{cold start}), ausência de parâmetros ou ausência em disco de arquivos auxiliares de metadados.
    \item \textbf{Diagrama Arquitetural de Fluxo:} Ilustra esquematicamente e de forma inequívoca a trajetória do dado da interface aos modelos e de volta à tela.
\end{enumerate}

% ==============================================================================
% 3. ETAPA 3: MAPEAMENTO DOS FLUXOS ATUAIS (FASES 1 E 2)
% ==============================================================================
\subsection{Etapa 3: Mapeamento dos Fluxos Atuais (Fases 1 e 2)}

Apresenta-se a seguir a especificação técnica dos fluxos implementados no sistema principal (\texttt{app/dashboard.py}). 

\subsubsection{Organização da Navegação Reativa e Modos de Visita}

O fluxo de abas atual tem por objetivo atender perfis e níveis de profundidade distintos, resultando na seguinte navegação:
\begin{itemize}
    \item \textbf{Comece Aqui (Apresentação Geral):} Visão panorâmica introdutória destinada a leigos e visitantes gerais;
    \item \textbf{Modo Candidato:} Abrange os fluxos analíticos e operacionais 1 a 3 descritos a seguir:
    \begin{enumerate}[label=\arabic*.]
        \item Visão de Negócio (EDA de Mercado);
        \item Simulador CBF (Filtragem Baseada em Conteúdo);
        \item Filtragem Colaborativa (SVD Latente).
    \end{enumerate}
    \item \textbf{Modo Avaliador:} Direcionado à professora da disciplina, incorporando:
    \begin{enumerate}[label=\arabic*., start=4]
        \item Comparativo CBF $\times$ CF (duelo formal com CBF métrica, oráculo/piso de ruído, Wilcoxon e proveniência);
        \item Como avaliamos \& limitações (Fluxo 5): declaração explícita de metodologia, dados sintéticos, conformidade com a LGPD, avaliação \textit{self-fulfilling}, ruído estocástico e nota do baseline de popularidade.
    \end{enumerate}
\end{itemize}

% ------------------------------------------------------------------------------
% 3.1 ABA 1: VISÃO DE NEGÓCIO (EDA)
% ------------------------------------------------------------------------------
\subsubsection{Aba 1: Visão de Negócio (EDA de Mercado)}

\noindent\textbf{User Story Vinculada:} US1 -- Visão de Mercado (\textit{Persona: Candidato}).
Declaração formal e critérios de aceite BDD na Seção~\ref{sec:us1_mercado}.

\begin{enumerate}[label=\textbf{\arabic*.}]
    \item \textbf{Objetivo:} Apresentar os achados empíricos do mercado de trabalho e submeter a validação formal as 5 hipóteses estratégicas de negócio, fornecendo embasamento quantitativo antes de se propor qualquer recomendação algorítmica.
    \item \textbf{Fontes de Dados \& Artefatos:}
    \begin{itemize}
        \item \texttt{data/raw/postings.csv}: Conjunto de dados primário contendo 123.849 anúncios de empregos;
        \item \texttt{data/raw/companies/companies.csv}: Cadastro de dados cadastrais e porte corporativo das empresas contratantes.
    \end{itemize}
    \item \textbf{Cadeia de Processamento (\texttt{src/gerar\_figuras\_eda.py}):}
    \begin{itemize}
        \item Higienização estatística, tratamento de registros incompletos e imputação de valores nulos;
        \item Agregação das taxas de conversão de mercado (\textit{Click-Through Rate} -- CTR), calculadas como a razão entre candidaturas registradas e visualizações de página:
        \begin{equation}
            \text{CTR} = \frac{\text{applies}}{\text{views}}
        \end{equation}
        \item Segmentação comparativa da mediana de CTR entre as modalidades presencial e remota;
        \item Mapeamento das faixas salariais por nível de senioridade e faixa de porte corporativo;
        \item Avaliação das 5 hipóteses estratégicas formuladas:
        \begin{itemize}
            \item \textbf{3 Confirmadas:} Maior preferência por vagas remotas ($H_1$); Disparidade salarial por senioridade ($H_2$); e Maior transparência salarial no remoto ($H_5$);
            \item \textbf{2 Refutadas:} Suposição de remuneração presencial mais alta como compensação de deslocamento ($H_3$); e Premissa de remuneração corporativa superior em empresas de grande porte ($H_4$).
        \end{itemize}
    \end{itemize}
    \item \textbf{Saída Visual:} Cartões diagnósticos de indicadores, relatórios de verificação das hipóteses estratégicas e gráficos de distribuição salarial e engajamento.
    \item \textbf{Comportamento sob Limitações / Fallback:} Opera de modo totalmente independente de identificação de usuário. Caso os dados brutos estejam inacessíveis, uma exceção controlada orienta a execução da rotina de carga via script.
\end{enumerate}

\begin{figure}[H]
\centering
\resizebox{0.92\textwidth}{!}{%
\begin{tikzpicture}[
    node distance=1.1cm,
    databox/.style={draw=black!70, rounded corners=3pt, fill=orange!10, font=\footnotesize, align=center, inner sep=6pt, text width=3.8cm},
    procbox/.style={draw=black!70, rounded corners=3pt, fill=blue!5, font=\footnotesize, align=center, inner sep=6pt, text width=4.5cm},
    calcbox/.style={draw=black!70, rounded corners=3pt, fill=green!8, font=\footnotesize, align=center, inner sep=6pt, text width=4.2cm},
    outbox/.style={draw=black!70, rounded corners=3pt, fill=purple!8, font=\footnotesize, align=center, inner sep=6pt, text width=3.8cm},
    arrow/.style={-{Stealth[length=2.5mm]}, thick, draw=black!75}
]

\node[databox] (raw_p) {\textbf{data/raw/postings.csv}\\(123.849 vagas)};
\node[databox, below=5mm of raw_p] (raw_c) {\textbf{data/raw/companies.csv}\\(segmentação empresarial)};

\node[procbox, right=3.2cm of $(raw_p)!0.5!(raw_c)$] (proc) {\textbf{Tratamento \& Limpeza}\\(\texttt{src/gerar\_figuras\_eda.py})\\Imputação e remoção de outliers};

\node[calcbox, right=1.4cm of proc] (calc) {\textbf{Cálculo de Indicadores}\\CTR, Faixa Salarial,\\Modalidade Remota vs. Presencial};

\node[outbox, right=1.4cm of calc] (ui) {\textbf{Aba 1 (Streamlit)}\\Cards de Hipóteses\\ e Gráficos};

\draw[arrow] (raw_p.east) -- (proc.west);
\draw[arrow] (raw_c.east) -- (proc.west);
\draw[arrow] (proc) -- (calc);
\draw[arrow] (calc) -- (ui);

\end{tikzpicture}%
}
\caption{Diagrama arquitetural de fluxo da Aba 1 (Visão de Negócio -- EDA).}
\label{fig:fluxo_aba1}
\end{figure}

% ------------------------------------------------------------------------------
% 3.2 ABA 2: SIMULADOR CBF
% ------------------------------------------------------------------------------
\subsubsection{Aba 2: Simulador CBF (Filtragem Baseada em Conteúdo)}

\noindent\textbf{User Story Vinculada:} US2 -- Filtragem por Conteúdo e Cold Start (\textit{Persona: Candidato}).
Declaração formal e critérios de aceite BDD na Seção~\ref{sec:us2_cold_start}.

\begin{enumerate}[label=\textbf{\arabic*.}]
    \item \textbf{Objetivo:} Permitir que o candidato construa em tempo real o seu perfil profissional sob demanda, declarando competências técnicas pontuais ou interagindo com vagas de interesse, eliminando a dependência de um histórico prévio de candidaturas.
    \item \textbf{Fontes de Dados \& Artefatos:}
    \begin{itemize}
        \item \texttt{data/raw/postings.csv}: Amostra representativa cacheada de 25.000 anúncios de vagas;
        \item \texttt{src/recomendador.py}: Módulo computacional com a classe central \texttt{RecSysCBF}.
    \end{itemize}
    \item \textbf{Cadeia de Processamento (Pipeline de Lógica):}
    Construção do identificador textual composto ponderado (título com peso dobrado $+$ competências $+$ senioridade), vetorização esparsa via TF-IDF (catálogo com 3.000 termos), síntese aditiva contínua do vetor de perfil do usuário $\mathbf{u}$ e mensuração por similaridade de cosseno combinada ao bônus empírico de CTR. O detalhamento analítico e as deduções matemáticas formais encontram-se na Seção~\ref{sec:fase1_cbf}.
    \item \textbf{Comportamento sob Limitações / Fallback (Cold Start):} Resolvido nativamente por construção. Por apoiar-se exclusivamente na semântica do item e nas competências digitadas, basta que o usuário informe uma única palavra-chave (e.g., ``Python'') para que o sistema gere recomendações altamente aderentes, dispensando qualquer histórico prévio.
\end{enumerate}

\begin{figure}[H]
\centering
\resizebox{0.95\textwidth}{!}{%
\begin{tikzpicture}[
    node distance=1.1cm,
    uibox/.style={draw=black!70, rounded corners=3pt, fill=orange!10, font=\footnotesize, align=center, inner sep=5pt, text width=3.8cm},
    procbox/.style={draw=black!70, rounded corners=3pt, fill=blue!5, font=\footnotesize, align=center, inner sep=5pt, text width=4.2cm},
    databox/.style={draw=black!70, rounded corners=3pt, fill=purple!8, font=\footnotesize, align=center, inner sep=5pt, text width=3.8cm},
    feedbox/.style={draw=black!70, rounded corners=3pt, fill=green!8, font=\footnotesize, align=center, inner sep=5pt, text width=4.2cm},
    statebox/.style={draw=black!70, rounded corners=3pt, fill=yellow!15, font=\footnotesize, align=center, inner sep=5pt, text width=3.8cm},
    arrow/.style={-{Stealth[length=2.5mm]}, thick, draw=black!75},
    farrow/.style={-{Stealth[length=2.5mm]}, thick, dashed, draw=blue!80}
]

% Linha Superior: Usuário -> Perfil -> Similaridade
\node[uibox] (ui_in) {\textbf{Usuário na Interface}\\Digita Skills ($\gamma$)\\Seleciona Likes ($\alpha$) / Dislikes ($\beta$)};
\node[procbox, right=1.6cm of ui_in] (build_prof) {\textbf{RecSysCBF.build\_user\_profile}\\(\texttt{src/recomendador.py})\\Vetor do Usuário $\mathbf{u} \in \mathbb{R}^{3000}$};
\node[procbox, right=1.6cm of build_prof] (cos_sim) {\textbf{Similaridade do Cosseno}\\$\cos(\mathbf{u}, \mathbf{d}_j) = \frac{\mathbf{u} \cdot \mathbf{d}_j}{\|\mathbf{u}\| \|\mathbf{d}_j\|}$};

% Matriz TF-IDF acima da Similaridade
\node[databox, above=1.0cm of cos_sim] (tfidf) {\textbf{Matriz TF-IDF Cacheada}\\(3.000 termos no \texttt{postings.csv})};

% Linha Intermediária: Score Final e Bônus CTR à direita
\node[procbox, below=1.2cm of cos_sim] (score_fin) {\textbf{Score Final \& Filtros}\\$\text{Score} = \cos + \alpha_{\text{CTR}}\cdot\text{CTR}$\\Corte Top-$N$ \& Filtro Remoto};
\node[databox, right=1.2cm of score_fin, fill=yellow!10] (ctr_bonus) {\textbf{Bônus Empírico CTR}\\$+ \alpha_{\text{CTR}} \cdot \text{CTR}_j$};

% Linha Inferior: Feed e Session State alinhados
\node[feedbox, below=2.4cm of build_prof] (feed) {\textbf{Feed de \texttt{card\_vaga}}\\Aderência semântica (\% match),\\Termos que casam, Feedback};
\node[statebox, left=2.0cm of feed] (state) {\textbf{session\_state}\\Listas de curtidas\\e descartadas};

% Conexões principais
\draw[arrow] (ui_in.east) -- (build_prof.west);
\draw[arrow] (build_prof.east) -- (cos_sim.west);
\draw[arrow] (tfidf.south) -- (cos_sim.north);
\draw[arrow] (cos_sim.south) -- (score_fin.north);
\draw[arrow] (ctr_bonus.west) -- (score_fin.east);
\draw[arrow] (score_fin.south) |- (feed.east);

% Loop de feedback interativo
\draw[farrow] (feed.west) -- node[above, font=\scriptsize] {Feedback Likes / Dislikes} (state.east);
\draw[farrow] (state.north) -- node[left, font=\scriptsize, align=center] {Re-ranqueia\\(\texttt{\_rerun\_feed()})} (ui_in.south);

\end{tikzpicture}%
}
\caption{Diagrama arquitetural de fluxo da Aba 2 (Simulador CBF com ciclo de feedback reativo).}
\label{fig:fluxo_aba2}
\end{figure}

% ------------------------------------------------------------------------------
% 3.3 ABA 3: FILTRAGEM COLABORATIVA (SVD)
% ------------------------------------------------------------------------------
\subsubsection{Aba 3: Filtragem Colaborativa}

\noindent\textbf{User Story Vinculada:} US3 -- Filtragem Colaborativa e Perfil Latente (\textit{Persona: Recrutador / Plataforma}).
Declaração formal e critérios de aceite BDD na Seção~\ref{sec:us_fase2}.

\begin{enumerate}[label=\textbf{\arabic*.}]
    \item \textbf{Objetivo:} Inspecionar o perfil latente implícito inferido a partir do comportamento do candidato e gerar recomendações de vagas baseadas no histórico coletivo de profissionais análogos, assegurando explicabilidade aditiva transparente.
    \item \textbf{Fontes de Dados \& Artefatos:}
    \begin{itemize}
        \item \texttt{modelo\_svd.pkl}: Modelo de Decomposição em Valores Singulares treinado e serializado;
        \item \texttt{data/processed/catalogo\_cf.csv}: Catálogo cacheado com 6.000 vagas;
        \item \texttt{data/processed/interacoes.parquet}: Matriz esparsa contendo 602.216 notas (escala de 1 a 5) atribuídas por 5.000 personas sintéticas estruturadas.
    \end{itemize}
    \item \textbf{Cadeia de Processamento (Pipeline de Lógica):}
    Recuperação do vetor de fatores latentes do usuário ($p_u \in \mathbb{R}^{20}$) e predição das notas estimadas via formulação aditiva clássica do SVD ($\hat{y}_{u,i} = \mu + b_u + b_i + q_i^T p_u$), seguida de ordenação decrescente Top-$N$ e decomposição interativa de explicabilidade (gráfico Waterfall Altair). O detalhamento analítico, os hiperparâmetros e a formulação matemática completa encontram-se na Seção~\ref{sec:fase2_cf}.
\end{enumerate}

\begin{figure}[H]
\centering
\resizebox{0.95\textwidth}{!}{%
\begin{tikzpicture}[
    node distance=1.1cm,
    uibox/.style={draw=black!70, rounded corners=3pt, fill=orange!10, font=\footnotesize, align=center, inner sep=5pt, text width=3.8cm},
    procbox/.style={draw=black!70, rounded corners=3pt, fill=blue!5, font=\footnotesize, align=center, inner sep=5pt, text width=4.4cm},
    databox/.style={draw=black!70, rounded corners=3pt, fill=purple!8, font=\footnotesize, align=center, inner sep=5pt, text width=4.0cm},
    feedbox/.style={draw=black!70, rounded corners=3pt, fill=green!8, font=\footnotesize, align=center, inner sep=5pt, text width=4.4cm},
    arrow/.style={-{Stealth[length=2.5mm]}, thick, draw=black!75}
]

\node[uibox] (u_sel) {\textbf{Seleção do Candidato}\\ID entre os 5.000 sintéticos};

\node[procbox, right=1.6cm of u_sel] (recup) {\textbf{Recuperação de Estado}\\Vetor latente $p_u \in \mathbb{R}^{20}$,\\Persona e Histórico de Treino};

\node[databox, above=8mm of recup] (models) {\textbf{Modelo Treinado}\\(\texttt{modelo\_svd.pkl})};

\node[databox, below=8mm of recup] (cat) {\textbf{Catálogo Cacheado}\\(\texttt{catalogo\_cf.csv} -- 6.000 vagas)};

\node[procbox, right=1.6cm of recup] (svd_pred) {\textbf{Predição SVD}\\$\hat{y} = \mu + b_u + b_i + q_i^T p_u$\\Ranking por Escore Bruto};

\node[feedbox, right=1.6cm of svd_pred] (waterfall) {\textbf{Waterfall Altair}\\$\mu \to +b_u \to +b_i \to +\text{match} \to =\hat{y}$\\Explicabilidade do Componente Dominante};

\node[feedbox, below=1.0cm of waterfall] (feed_cf) {\textbf{Feed Top-$N$ (\texttt{card\_vaga})}\\Exclui avaliadas, nota em destaque,\\Simulação didática de Curtir/Descartar};

\draw[arrow] (u_sel.east) -- (recup.west);
\draw[arrow] (models.south) -- (recup.north);
\draw[arrow] (cat.north) -- (recup.south);
\draw[arrow] (recup.east) -- (svd_pred.west);
\draw[arrow] (svd_pred.east) -- (waterfall.west);
\draw[arrow] (svd_pred.south) |- (feed_cf.west);

\end{tikzpicture}%
}
\caption{Diagrama arquitetural de fluxo da Aba 3 (Filtragem Colaborativa SVD com explicabilidade aditiva).}
\label{fig:fluxo_aba3}
\end{figure}

% ------------------------------------------------------------------------------
% 3.4 ABA 4: COMPARATIVO CBF X CF
% ------------------------------------------------------------------------------
\subsubsection{Aba 4: Comparativo CBF $\times$ CF (Estado Atual nas Fases 1 e 2)}

\noindent\textbf{User Story Vinculada:} US5 -- Auditoria e Duelo de Modelos (\textit{Persona: Professora / Avaliador}).
Declaração formal e critérios de aceite BDD na Seção~\ref{sec:us5_avaliador}.

\begin{enumerate}[label=\textbf{\arabic*.}]
    \item \textbf{Objetivo:} Estabelecer uma comparação analítica, teórica e metrológica rigorosa entre as abordagens de Filtragem Baseada em Conteúdo e Filtragem Colaborativa, fundamentando decisões de arquitetura de software sob evidências empíricas.
    \item \textbf{Fontes de Dados \& Artefatos:}
    \begin{itemize}
        \item \texttt{data/processed/metadados\_cf.pkl}: Dicionário gerado pelo script oficial de benchmarking (\texttt{src/executar\_modelagem.py}), integrando as métricas de SVD, KNN Cosseno, baselines triviais, piso do oráculo e p-valor do teste de Wilcoxon;
        \item \texttt{data/processed/metadados\_cbf.pkl}: Metadados de avaliação do CBF sob idêntico protocolo experimental, gerados por \texttt{src/avaliar\_cbf.py};
        \item Carga e unificação gerenciadas pela função centralizadora \texttt{carregar\_metricas()}.
    \end{itemize}
    \item \textbf{Cadeia de Processamento (Pipeline de Lógica):}
    \begin{itemize}
        \item Estruturação da tabela auditada com 8 linhas de modelos:
        \begin{enumerate}[label=\roman*.]
            \item Chute Trivial 1: Média Global ($\mu$);
            \item Chute Trivial 2: Média por Item ($b_i$);
            \item Chute Trivial 3: Aleatório Uniforme;
            \item Baseline de Popularidade: ranking por mediana de CTR do treino em pool de 2.000 itens;
            \item Oráculo Teórico (Piso Irredutível de Ruído): simulação do limite bayesiano de erro com ruído gaussiano $\sigma = 0{,}520$;
            \item CBF Métrica: avaliação de Precision@10 e NDCG@10 na base de teste sem a aplicação do bônus de CTR, garantindo isonomia com a formulação pura;
            \item KNN Cosseno Centrado (\textit{User-based Collaborative Filtering});
            \item SVD Fatores Latentes ($k=20$).
        \end{enumerate}
        \item Apresentação das métricas de ranqueamento (Precision@10, NDCG@10 e RMSE) acompanhadas do p-valor do teste não-paramétrico de Wilcoxon, atestando significância estatística da superioridade do SVD;
        \item Registro transparente da data e horário de geração dos artefatos (\textit{proveniência de dados});
        \item Confrontação conceitual multidimensional (\textit{Sinal Primário}, \textit{Perfil do Usuário}, \textit{Vulnerabilidade a Cold Start}, \textit{Potencial de Serendipidade} e \textit{Explicabilidade}).
    \end{itemize}
    \item \textbf{Comportamento sob Limitações / Fallback:} Se o artefato \texttt{metadados\_cbf.pkl} não for localizado em disco, a renderização da tabela adota tratamento defensivo através de \texttt{.get()}, exibindo o marcador textual ``---'' nas linhas de conteúdo sem interromper o funcionamento ou causar falhas na página.
\end{enumerate}

\begin{figure}[H]
\centering
\resizebox{0.92\textwidth}{!}{%
\begin{tikzpicture}[
    node distance=1.1cm,
    databox/.style={draw=black!70, rounded corners=3pt, fill=orange!10, font=\footnotesize, align=center, inner sep=6pt, text width=4.2cm},
    procbox/.style={draw=black!70, rounded corners=3pt, fill=blue!5, font=\footnotesize, align=center, inner sep=6pt, text width=4.4cm},
    outbox/.style={draw=black!70, rounded corners=3pt, fill=green!8, font=\footnotesize, align=center, inner sep=6pt, text width=4.4cm},
    arrow/.style={-{Stealth[length=2.5mm]}, thick, draw=black!75}
]

\node[databox] (mcf) {\textbf{executar\_modelagem.py}\\$\to$ \texttt{metadados\_cf.pkl}\\(SVD, KNN, Baselines, Oráculo)};
\node[databox, below=8mm of mcf] (mcbf) {\textbf{avaliar\_cbf.py}\\$\to$ \texttt{metadados\_cbf.pkl}\\(usa \texttt{protocolo\_ranking.pkl})};

% Ponto médio seguro na borda direita das caixas de entrada
\path (mcf.east) -- (mcbf.east) coordinate[midway] (east_mid);
\node[procbox, anchor=west, right=1.6cm of east_mid] (merge) {\textbf{carregar\_metricas()}\\Unificação e tratamento defensivo\\Fallback automático com ``---''};

\node[outbox, right=1.6cm of merge] (duelo) {\textbf{Duelo dos Modelos (UI)}\\Tabela 8 linhas, Wilcoxon,\\Proveniência \& Matriz Conceitual};

\draw[arrow] (mcf.east) -- (merge.west);
\draw[arrow] (mcbf.east) -- (merge.west);
\draw[arrow] (merge) -- (duelo);

\end{tikzpicture}%
}
\caption{Diagrama arquitetural de fluxo da Aba 4 (Comparativo CBF $\times$ CF).}
\label{fig:fluxo_aba4}
\end{figure}

% ------------------------------------------------------------------------------
% 3.5 FLUXO 5: COMO AVALIAMOS & LIMITAÇÕES
% ------------------------------------------------------------------------------
\subsubsection{Fluxo 5: Como avaliamos \& limitações (Modo Avaliador / Professora)}

\noindent\textbf{User Story Vinculada:} US5 -- Transparência, Governança e Limitações (\textit{Persona: Professora / Avaliador}).
Declaração formal e critérios de aceite BDD na Seção~\ref{sec:us5_avaliador}.

\begin{enumerate}[label=\textbf{\arabic*.}]
    \item \textbf{Objetivo:} Expor a metodologia, o processo de validação e as fragilidades do sistema explicitamente na interface visual, a fim de superar a ocultação da tomada de decisão e limitações em arquivos de texto auxiliares.
    \item \textbf{Fontes de Dados \& Artefatos:}
    \begin{itemize}
        \item Metadados de timestamp de modificação (\textit{mtime}) dos artefatos em \texttt{data/processed/};
        \item Chaves serializadas em \texttt{metadados\_cf.pkl} e \texttt{metadados\_cbf.pkl} (\texttt{rmse\_oraculo}, \texttt{p\_wilcoxon}, \texttt{precision\_k}, \texttt{ndcg\_k}).
    \end{itemize}
    \item \textbf{Conteúdo Estruturado}
    \begin{itemize}
        \item \textbf{Tabela de Proveniência de Artefatos:} Mapeamento inequívoco de cada artefato binário gerado, indicando o script responsável pela criação, a base de dados alimentadora e a data exata da compilação;
        \item \textbf{Conformidade com a LGPD:} Fundamentação jurídica e técnica do uso de dados de interação sintéticos estruturados (Lei Federal nº 13.709/2018), resguardando a privacidade e a confidencialidade de profissionais reais do mercado;
        \item \textbf{Avaliação \textit{Self-Fulfilling}:} Reconhecimento epistêmico de que a geração do \textit{ground truth} mediante decomposição bilinear impõe uma circularidade favorável ao modelo SVD, exigindo a mensuração da proximidade com o oráculo para interpretação fidedigna;
        \item \textbf{Piso de Ruído Estocástico:} Comparativo entre o erro obtido pelo SVD ($0{,}625$) e o limite bayesiano do oráculo ($0{,}521$), atestando uma razão de apenas $1{,}20\times$ em relação ao limite estocástico ideal;
        \item \textbf{Limitações Abertas:} Registro público de fragilidades do catálogo primário, incluindo a esparsidade do campo \texttt{skills\_desc} ($\approx 98\%$ de ausências), a degradação e falha do algoritmo KNN em matrizes ultrarralas e a dependência do motor híbrido para mitigação do \textit{cold start} na Fase 3.
    \end{itemize}
    \item \textbf{Fallback e Tratamento Defensivo:} Cada bloco de conteúdo analítico é resguardado pela inspeção da presença das respectivas chaves no dicionário de metadados; na eventualidade de execução com arquivos de métricas legados, o componente omite os blocos não preenchidos sem degradar o restante da página.
\end{enumerate}

% ==============================================================================
% 4. ETAPA 4: EVOLUÇÃO PARA A FILTRAGEM HÍBRIDA (FASE 3)
% ==============================================================================
\subsection{Etapa 4: Evolução para a Filtragem Híbrida (Fase 3)}

O projeto passará a implementar futuramente a filtragem híbrida, criando assim uma nova aba para análise dessa abordagem.

\subsubsection{Estratégias de Hibridização (Taxonomia de Robin Burke)}

Para neutralizar a propensão à hipersensibilidade do CBF e a vulnerabilidade ao \textit{cold start} do CF, adota-se a conjugação de duas técnicas clássicas propostas por Robin Burke \cite{burke2002}: fusão linear ponderada (\textit{Weighted}, com pesos empíricos $w_1 = 0{,}40$ e $w_2 = 0{,}60$) sobre escores previamente normalizados e chaveamento dinâmico de regime (\textit{Switching}) para segregar novos usuários ($|I_u| = 0$). O equacionamento matemático formal e a formulação analítica completa encontram-se detalhados na Seção~\ref{sec:fase3_hibrido}.

\subsubsection{Diagrama do Fluxo da Arquitetura Alvo Híbrida (Fase 3)}

A Figura~\ref{fig:fluxo_hibrido_fase3} documenta a arquitetura de dados e de tomada de decisão projetada para o motor híbrido, decomposta em seus 5 subgrafos computacionais.

\begin{figure}[H]
\centering
\resizebox{0.95\textwidth}{!}{%
\begin{tikzpicture}[
    node distance=1.2cm,
    box/.style={draw=black!70, rounded corners=3pt, fill=white, font=\footnotesize, align=center, inner sep=6pt},
    decbox/.style={diamond, draw=black!70, fill=orange!15, font=\footnotesize, align=center, inner sep=2pt, aspect=1.8},
    subg/.style={draw=black!40, dashed, rounded corners=6pt, fill=black!2, inner sep=8pt},
    arrow/.style={-{Stealth[length=2.5mm]}, thick, draw=black!75},
    darrow/.style={-{Stealth[length=2.5mm]}, dashed, thick, draw=blue!80}
]

% Subgrafo 1: Entrada do Usuário
\node[box, fill=gray!10, text width=5.0cm] (u_node) {\textbf{Usuário na Interface}\\(Identificado com ID ou Convidado)};

\begin{scope}[on background layer]
    \node[subg, fit=(u_node), label={[font=\small\bfseries]above:1. Entrada do Usuário}] (sg_in) {};
\end{scope}

% Subgrafo 2: Decisão de Cold Start (Switching)
\node[decbox, below=1.3cm of u_node, text width=3.4cm] (cond_cold) {Possui histórico\\no sistema?\\($|I_u| > 0$)};

\begin{scope}[on background layer]
    \node[subg, fit=(cond_cold), label={[font=\small\bfseries]above:2. Decisão de Cold Start (Switching)}] (sg_dec) {};
\end{scope}

% Subgrafo 3: Motores de Recomendação (Lado a Lado)
\node[box, below left=1.6cm and 0.8cm of cond_cold, text width=4.4cm, fill=blue!8] (engine_cbf) {\textbf{Motor CBF (TF-IDF)}\\$\text{Score}_{\text{CBF}}^{\text{norm}} \in [0, 1]$\\Similaridade textual e competências};
\node[box, below right=1.6cm and 0.8cm of cond_cold, text width=4.4cm, fill=blue!8] (engine_cf) {\textbf{Motor CF (SVD Normalizado)}\\$\text{Score}_{\text{CF}}^{\text{norm}} \in [0, 1]$\\Fatores latentes e comportamento};

\begin{scope}[on background layer]
    \node[subg, fit=(engine_cbf) (engine_cf), label={[font=\small\bfseries]above:3. Motores de Recomendação}] (sg_eng) {};
\end{scope}

% Subgrafo 4: Fusão & Ranking (Claramente abaixo dos motores)
\path (engine_cbf.south) -- (engine_cf.south) coordinate[midway] (eng_south);
\node[box, below=1.5cm of eng_south, text width=7.0cm, fill=green!10] (fusao) {\textbf{Fusão Ponderada (Weighted)}\\$\text{Score}_{\text{Final}} = w_1 \cdot \text{Score}_{\text{CBF}} + w_2 \cdot \text{Score}_{\text{CF}}$\\(Ex.: $0{,}40\text{ CBF} + 0{,}60\text{ CF}$)};

\node[box, below=1.0cm of fusao, text width=5.2cm, fill=purple!8] (rank) {\textbf{Top-$N$ Recomendações Híbridas}\\Desempate por CTR e regras de negócio};

\begin{scope}[on background layer]
    \node[subg, fit=(fusao) (rank), label={[font=\small\bfseries]above:4. Fusão \& Ranking}] (sg_fuse) {};
\end{scope}

% Subgrafo 5: Avaliação Comparativa
\node[box, below=1.4cm of rank, text width=7.4cm, fill=yellow!15] (comp) {\textbf{Painel Comparativo e Auditoria}\\Métricas Isoladas vs. Ganhos da Hibridização\\(Mitigação de Cold Start + Serendipidade)};

\begin{scope}[on background layer]
    \node[subg, fit=(comp), label={[font=\small\bfseries]above:5. Avaliação Comparativa}] (sg_comp) {};
\end{scope}

% Conexões
\draw[arrow] (u_node) -- (cond_cold);

% Decisão de chaveamento (Saídas para direita e esquerda sem sobreposição)
\draw[arrow] (cond_cold.east) -| node[above, font=\scriptsize, pos=0.35] {Sim ($|I_u| > 0$)} (engine_cf.north);
\draw[arrow] (cond_cold.west) -| node[above, font=\scriptsize, pos=0.35] {Não (Cold Start)} (engine_cbf.north);

% Alimentação da Fusão Ponderada
\draw[arrow] (engine_cbf.south) |- (fusao.west);
\draw[arrow] (engine_cf.south) |- (fusao.east);

% Rota de Chaveamento para Cold Start (contorna pela esquerda diretamente ao ranking)
\draw[darrow] (engine_cbf.west) -- ++(-1.0, 0) |- node[pos=0.75, left, font=\scriptsize, text=blue!80] {Chaveamento (100\% CBF)} (rank.west);

% Encadeamento final
\draw[arrow] (fusao) -- (rank);
\draw[arrow] (rank) -- (comp);

\end{tikzpicture}%
}
\caption{Diagrama arquitetural completo da Filtragem Híbrida (Fase 3), combinando chaveamento de cold start (\textit{Switching}) com combinação linear normalizada (\textit{Weighted}).}
\label{fig:fluxo_hibrido_fase3}
\end{figure}
